% Chapter 22 -- Claim: Transformers Price Stocks Better Than Factor Models
% Requires: amsmath, amssymb, booktabs, array, tabularx, longtable, tcolorbox (breakable, skins), hyperref or url.
% Defines the callout box only if the preamble has not.
\providecolor{fenavy}{HTML}{1B2A6B}
\makeatletter
\@ifundefined{callout}{%
  \newtcolorbox{callout}{breakable, enhanced,
    colback=white, colframe=gray!70, boxrule=0.5pt, arc=1mm,
    borderline west={2.5pt}{0pt}{gray!70},
    left=8pt, right=8pt, top=4pt, bottom=4pt}%
}{}
\makeatother

\providecolor{fenavy}{HTML}{1B2A6B}
\providecolor{feshade}{HTML}{F3F5FA}
\makeatletter
% Shaded box for optional technical material that general readers may skip.
\@ifundefined{technote}{%
  \newtcolorbox{technote}{breakable, enhanced,
    colback=feshade, colframe=fenavy!60, boxrule=0.5pt, arc=1mm,
    left=8pt, right=8pt, top=4pt, bottom=4pt}%
}{}
\makeatother

\chapter{Claim: Transformers Price Stocks Better Than Factor Models}
\label{ch:22}
\providecommand{\term}[1]{\emph{#1}}

Anyone who follows the stock market notices that companies do not move alone. When one bank reports trouble, other banks' shares often fall too, and a supplier's bad news can drag down its customers. The newest artificial-intelligence models are designed to learn such connections from data. This chapter asks whether letting a model learn how stocks relate to one another really explains their returns better than the factor models of earlier chapters.

Asset pricing asks why different assets earn different average returns. Since the 1980s its general answer has been the \term{stochastic discount factor}\index{stochastic discount factor (SDF)}, a single random variable that prices every asset, which can be thought of as a gauge of how bad economic times are, used to weigh every payoff. The factor models of Chapter~\ref{ch:03} are one way to build it, from a handful of portfolios chosen by researchers. With hundreds of candidate characteristics, later work asked the data to build it instead. Kozak, Nagel and Santosh (2020) estimated a discount factor from many anomaly portfolios with shrinkage, which keeps so many portfolios from fitting noise. Kelly, Pruitt and Su (2019), noting that a stock's risk changes with its characteristics, let each stock's exposures depend on them. Chen, Pelger and Zhu (2024) used deep neural networks to allow more flexible relationships. Each step priced returns more accurately than the last.

The newest step comes from the technology behind large language models. A \textbf{transformer}\index{transformer|textbf} uses a mechanism called \term{attention}\index{attention (neural networks)|textbf}, which lets each word in a sentence look at the other words to decide what it means. Kelly, Kuznetsov, Malamud and Xu (2025) placed a transformer inside an asset-pricing model, so that each stock could ``look at'' other stocks before the model decided how to price it. They reported that this artificial intelligence pricing model priced U.S. stocks with errors about 30\% smaller than the best earlier models. Its tangency portfolio, the portfolio with the highest Sharpe ratio, earned an out-of-sample Sharpe ratio of about four. That is far above anything an investor has earned for long, so the result calls for careful testing.

The stakes are large if the result holds. It would mean that the price of a stock depends on information about other stocks in ways that factor models cannot capture. It would also mean that the methods behind chatbots can improve one of the central models of finance. 

The claim on trial\index{claims-on-trial method} is: \textbf{an asset-pricing model that places a transformer inside the stochastic discount factor prices the cross-section with much smaller errors than earlier models, and its tangency portfolio earns a higher Sharpe ratio out of sample.}

The lab in Section~\ref{sec:22-lab} replicates the linear-transformer pricing model of Kelly et al.\ (2025) on free anomaly portfolios from the Jensen--Kelly--Pedersen data. It compares the model with the shrinkage baseline of Kozak, Nagel and Santosh (2020). The goal is to check whether attention across assets really prices returns better than a well-built factor model. The original study uses stock-level data from paid databases, so the lab replicates the mechanism one level up, with portfolios as the assets.


\section{Foundations}
\label{sec:22-foundations}

\subsection{The SDF and the pricing equation for excess returns}

An asset-pricing model has one job: to explain why some assets earn higher average returns than others. The most general way to write such a model uses a \textbf{stochastic discount factor}\index{stochastic discount factor|textbf} (SDF). This is a random number $M_{t+1}$ that is high in bad times and low in good times. The model says that every excess return\index{return!excess} $R^e_{i,t+1}$, once multiplied by the SDF, has an expected value of zero:
\begin{equation}
  E_t\!\left[ M_{t+1} R^e_{i,t+1} \right] = 0 \quad \text{for every asset } i.
  \label{eq:22-pricing}
\end{equation}
The intuition is simple. An asset that pays off in bad times, when $M$ is high, is valuable insurance, so investors accept a low average return on it. An asset that pays off in good times must offer a high average return. The CAPM of Chapter~\ref{ch:03} is a special case, in which $M$ falls in a straight line as the market return rises.

Many modern studies write the SDF as one minus the return on a portfolio of the assets:
\begin{equation}
  M_{t+1} = 1 - w_t' R^e_{t+1},
  \label{eq:22-sdf}
\end{equation}
Here $R^e_{t+1}$ is the vector of excess returns, and $w_t$ is the vector of portfolio weights, chosen with information available at time $t$. Estimating an asset-pricing model then means choosing the weights $w_t$. If the weights may depend on each stock's characteristics, such as its size, value and past return, the model can price the cross-section flexibly.

\subsection{From an SDF to the tangency portfolio}

In plain words, finding the discount factor turns out to be the same as finding the portfolio with the best possible Sharpe ratio, the tangency portfolio of Chapter~\ref{ch:09}. The shaded box shows why.

\begin{technote}
\textbf{Technical note: from the SDF to the tangency portfolio.} When the weights are constant, Equation~\eqref{eq:22-pricing} has a familiar solution. Substituting Equation~\eqref{eq:22-sdf} and solving for $w$ gives
\begin{equation}
  w = \Sigma^{-1} \mu,
  \label{eq:22-tangency}
\end{equation}
Here $\mu$ is the vector of expected excess returns and $\Sigma$ is the covariance matrix. (Strictly, second moments replace covariances, a small difference for monthly returns.) These are the weights of the tangency portfolio\index{tangency portfolio} of Chapter~\ref{ch:09}. The SDF portfolio is the portfolio with the highest possible Sharpe ratio, and its squared Sharpe ratio is $\mu'\Sigma^{-1}\mu$.
\end{technote}

This link is why studies of machine-learning SDFs report Sharpe ratios. A model that prices all assets correctly must contain the maximum-Sharpe portfolio. A model that leaves pricing errors behind has missed some return that a better portfolio would capture. Small pricing errors and a high out-of-sample Sharpe ratio are therefore two views of the same success. Studies measure pricing errors in two ways. One is the alphas\index{alpha} of test assets against the SDF portfolio. The other is the \term{Hansen--Jagannathan distance}\index{Hansen--Jagannathan distance|textbf}, which summarizes how far the model's SDF is from pricing a set of test assets exactly.

\subsection{Baseline 1: the shrinkage SDF of Kozak, Nagel and Santosh}

With hundreds of candidate factors, Equation~\eqref{eq:22-tangency} cannot be used directly. The sample $\Sigma$ is noisy, its inverse is noisier still, and the weights become extreme. Chapter~\ref{ch:09} called this problem error maximization. Kozak, Nagel and Santosh (2020) shrink the weights with a ridge penalty, exactly as in Chapter~\ref{ch:20}:
\begin{equation}
  \hat w = \left(\hat\Sigma + \gamma I\right)^{-1} \hat\mu,
  \label{eq:22-kns}
\end{equation}
Here $\gamma$ controls how strongly the weights are pulled toward zero, and it is chosen by cross-validation. The penalty is easiest to understand through the \term{principal components}\index{principal components} of the factors. It shrinks the weight on each component by the factor $\lambda_j/(\lambda_j+\gamma)$, where $\lambda_j$ is the component's variance. Components with little variance are shrunk the most, because a large average return on them is most likely to be noise. There is also an economic reason. A factor with a high average return but almost no variance would be close to an arbitrage opportunity, and such opportunities should be rare.

Kozak, Nagel and Santosh found that an SDF built from a few characteristic-based factors, such as the four- and five-factor models, could not summarize the cross-section of expected returns well, while an SDF built from a small number of principal components of many factors performed well.

\subsection{Baseline 2: IPCA}

\textbf{Instrumented principal components analysis}\index{IPCA|textbf} (IPCA), introduced by Kelly, Pruitt and Su (2019), lets factor loadings change with a stock's characteristics. The model is
\begin{equation}
  r_{i,t+1} = \beta_{i,t}' f_{t+1} + \varepsilon_{i,t+1}, \qquad \beta_{i,t} = \Gamma' z_{i,t},
  \label{eq:22-ipca}
\end{equation}
Here $f_{t+1}$ is a short vector of unobserved factors, $z_{i,t}$ is the stock's characteristics, and $\Gamma$ is a matrix estimated from the data. A stock's characteristics therefore tell the model how much risk it carries, which is the meaning of the paper's title, ``Characteristics Are Covariances.'' Kelly, Pruitt and Su found that five IPCA factors explained the cross-section of average returns more accurately than existing factor models. The remaining anomaly alphas were small and insignificant, and only about ten characteristics mattered.

\subsection{Deep SDFs: feed-forward networks and adversarial estimation}

A natural next step lets a neural network choose the weights in Equation~\eqref{eq:22-sdf}: $w_{i,t} = g(z_{i,t}, s_t)$. Here $g$ is a network of the kind described in Section~\ref{sec:20-nn}, and $s_t$ is a summary of economic conditions. The network is trained to make the pricing errors as small as possible. This is the same as making the SDF portfolio's Sharpe ratio as large as possible.

Chen, Pelger and Zhu (2024) added a twist from machine learning called \term{adversarial training}\index{adversarial training|textbf}. A second network searches for the portfolios that the SDF prices worst. The first network then improves its pricing of exactly those portfolios. The contest continues until the adversary can find no badly mispriced portfolio. They used U.S. stocks from 1967 to 2016, 46 firm characteristics and 178 macroeconomic series, which they summarized with a recurrent network. Their SDF portfolio reached an annual out-of-sample Sharpe ratio of about 2.6, roughly twice that of a network trained only to forecast returns.

\subsection{Attention and the transformer for a panel of stocks}

The models above treat each stock separately: stock $i$'s weight depends only on stock $i$'s own characteristics. Attention removes that restriction. In plain words, each stock looks at every other stock, judges how similar they are, and borrows information from the most similar ones. Give each stock a vector of characteristics $x_i$. Attention computes three transformed versions of it: a \term{query} $q_i = W_Q x_i$, a \textbf{key} $k_i = W_K x_i$ and a \textbf{value} $v_i = W_V x_i$. Stock $i$ compares its query with every stock's key, which gives a score for each stock. The scores are turned into weights that sum to one, and stock $i$ takes the weighted average of the other stocks' values:
\begin{equation}
  a_{ij} = \frac{\exp\left(q_i' k_j/\sqrt{d}\right)}{\sum_{l} \exp\left(q_i' k_l/\sqrt{d}\right)}, \qquad h_i = \sum_{j} a_{ij}\, v_j,
  \label{eq:22-attention}
\end{equation}
where $d$ is the length of the query vector. The output $h_i$ is a new description of stock $i$ that includes information from the stocks most similar to it, in the sense that the model has learned.

\textbf{Worked example.} Three stocks have queries and keys of length one, $q = (1, 0, 2)$ and $k = (1, 0, 2)$, and values $v = (0.03, 0.00, -0.02)$. For stock 1, the scores $q_1 k_j$ are $1, 0, 2$. The attention weights are $e^1, e^0, e^2$ divided by their sum, about $0.24, 0.09, 0.67$. Stock 1's new value is $0.24 \times 0.03 + 0.09 \times 0 + 0.67 \times (-0.02) \approx -0.006$. Stock 1 pays most attention to stock 3, whose key is most aligned with its query.

A transformer stacks many such attention steps; the shaded box shows the simplest version used in asset pricing.

\begin{technote}
\textbf{Technical note: the linear transformer.} A \textbf{transformer}\index{transformer|textbf} stacks several such attention layers, each followed by a small feed-forward network, with shortcuts that add each layer's input to its output. A \term{linear transformer} drops the exponential and the normalization. In the asset-pricing model of Kelly and coauthors, the linear version gives portfolio weights
\begin{equation}
  w_t = \left(X_t W X_t'\right) X_t \lambda,
  \label{eq:22-linear}
\end{equation}
where $X_t$ stacks all stocks' characteristics in rows, $W$ is a matrix of attention parameters and $\lambda$ a vector of pricing coefficients. The matrix $X_t W X_t'$ has one row and one column for every stock, and its entry $(i,j)$ says how much stock $i$ borrows from stock $j$.
\end{technote}

\section{How it works}
\label{sec:22-how}

\textbf{Why sharing information across stocks might help.} A stock's own characteristics are not the only information about its risk and return. News about a supplier affects its customers. A rise in one bank's funding costs affects other banks. Two stocks with similar characteristics may respond to the same shocks. Earlier models could use such links only if a researcher wrote them in by hand, for example through industry factors. Attention lets the model learn from the data, month by month, which other stocks are informative. Kelly and coauthors checked familiar links: shared analyst coverage, customer--supplier relationships and similar patents. These explained only about 18\% of the variation in their transformer's attention, so most of what the model learned lay outside those links.

\textbf{What the linear surrogate reveals.} Equation~\eqref{eq:22-linear} looks complicated, but it can be rewritten as an ordinary linear model in a very large set of features: every product of three characteristics. With 132 characteristics there are about 2.3 million such products. The linear transformer is therefore a ridge regression of the kind used throughout Chapter~\ref{ch:20}, on an enormous number of interactions. Its weights can be read as signals for timing portfolios built on characteristics. This rewriting is what makes the model, and this chapter's lab, understandable.

\textbf{More parameters than months.} The full transformer has around a million parameters, estimated in rolling windows of 60 months. Classical statistics would expect such a model to overfit. Chapter~\ref{ch:20} described the ``virtue of complexity'' argument, in which very large models with shrinkage can still forecast well. It also described Nagel's (2025) warning that such models may simply reproduce a simple strategy. Kelly, Kuznetsov, Malamud, Xu and Zhang (2024), in ``Large and Deep Factor Models,'' studied deep networks trained to build an SDF. They showed that such a network splits into two parts. One part discovers useful nonlinear characteristics. The other part prices them and behaves like a linear factor model with ridge shrinkage. Their analysis gives a theoretical reason why very large networks need not overfit, and a way to see what they have learned.

\section{The claim}
\label{sec:22-claim}

Stated so that it can fail:

\begin{callout}
\textbf{An asset-pricing model that uses a transformer to share information across stocks prices the cross-section of U.S. stock returns with smaller out-of-sample pricing errors than factor models and earlier machine-learning models, and the portfolio it implies earns a higher out-of-sample Sharpe ratio, among large stocks and after trading costs.}
\end{callout}

The \textbf{statistical form} asks for smaller pricing errors and a higher Sharpe ratio before costs; the \textbf{economic form} requires the advantage to survive among stocks an investor can trade in size and after costs. The tests are whether the gains appear \textbf{out of sample}, including after the early decades; whether the transformer beats both simple factor models and strong machine-learning \textbf{benchmarks}; whether the gains hold among mega and large stocks (\textbf{size}); and how much remains after \textbf{costs}.

\section{The original evidence}
\label{sec:22-evidence}

\subsection{Chen, Pelger and Zhu (2024)}

Chen, Pelger and Zhu\index{Chen, Pelger and Zhu} estimated an SDF for individual stocks with deep learning, training the network on the pricing condition of Equation~\eqref{eq:22-pricing} and using an adversary to find the hardest test portfolios. Their SDF beat their benchmarks in Sharpe ratio, explained variation and pricing errors, with familiar characteristics such as past price trends, profitability and capital structure doing most of the work, and with the economic environment also mattering.

\subsection{Kelly, Kuznetsov, Malamud, Xu and Zhang (2024)}

This working paper, ``Large and Deep Factor Models,'' shows that a deep SDF network works in two stages: it discovers nonlinear characteristics, then prices them like a large linear factor model. An earlier version reported that performance improved with depth. The paper has been substantially revised, so its empirical results should be read in their latest version.

\subsection{Kelly, Kuznetsov, Malamud and Xu (2025)}

Kelly, Kuznetsov, Malamud and Xu\index{AI pricing model (AIPM)|textbf} used monthly U.S. returns from 1963 to 2022 from the open data of Jensen, Kelly and Pedersen (2023), also used in Chapter~\ref{ch:16}, with 132 characteristics, the tiniest stocks dropped and the 132 matching factors as test assets. Each model was estimated in rolling 60-month windows, giving out-of-sample returns from 1968 to 2022, and the neural models were averaged over ten random starts (Table~\ref{tab:22-aipm}).

\begingroup
\small
\setlength{\LTpre}{6pt}\setlength{\LTpost}{6pt}
\begin{longtable}{@{}lcc@{}}
\caption{Out-of-sample performance of SDF models for U.S. stocks, 1968--2022, from Kelly, Kuznetsov, Malamud and Xu (2025), Table 1. The Sharpe ratio is annualized. The pricing error is the Hansen--Jagannathan distance on the 132 test factors; lower is better.}
\label{tab:22-aipm} \\
\toprule
Model & Sharpe ratio & Pricing error \\
\midrule
\endfirsthead
\multicolumn{3}{@{}l}{\emph{(continued)}} \\
\toprule
Model & Sharpe ratio & Pricing error \\
\midrule
\endhead
\midrule
\multicolumn{3}{r@{}}{\emph{(continued on the next page)}} \\
\endfoot
\bottomrule
\endlastfoot
Fama--French six factors & 1.01 & 0.55 \\*
Hou--Xue--Zhang $q$-factors & 1.77 & 0.42 \\
Linear in characteristics (no attention) & 3.60 & 0.15 \\
Shallow random-feature network & 3.87 & 0.13 \\
Linear attention & 3.89 & 0.14 \\
Deep network without attention (MLP) & 4.31 & 0.13 \\*
Transformer (AIPM) & 4.57 & 0.09 \\
\end{longtable}
\endgroup

The transformer had the highest Sharpe ratio and the smallest pricing error, about 30\% below the deep network without attention, with an alpha against it of about 1.9\% a year ($t = 8.4$). Its advantage held after 2002, with a Sharpe ratio of 3.37 against 3.07. More attention blocks helped, raising the Sharpe ratio from about 3.8 with one to about 4.6 with ten, while larger networks without attention did not. And it was best within each size group, though its Sharpe ratio was much lower among mega and large stocks than among small and micro stocks.

\subsection{Questions the papers leave open}

\begin{itemize}
  \item \textbf{Costs.} No trading costs are reported, and portfolios re-estimated monthly over thousands of stocks can trade heavily (Chapter~\ref{ch:14}).
  \item \textbf{Small stocks.} The highest Sharpe ratios come from small and micro stocks, where Chapters~\ref{ch:16} and~\ref{ch:20} found anomalies and machine-learning gains concentrated.
  \item \textbf{Replication.} The result comes from one team's working paper and has not been reproduced independently. Every model in Table~\ref{tab:22-aipm} reports a Sharpe ratio far above what investors have earned, a reminder that these returns are hypothetical and frictionless.
\end{itemize}

\section{Lab~\thechapter: Replicating the linear-transformer SDF on anomaly portfolios}
\label{sec:22-lab}

\textbf{Claim on trial\index{claims-on-trial method}:} attention across assets improves an SDF's out-of-sample pricing and Sharpe ratio. The lab replicates Kelly, Kuznetsov, Malamud and Xu (2025) against the baseline of Kozak, Nagel and Santosh (2020). Stock-level characteristics require paid CRSP and Compustat data, so the lab replicates the mechanism one level up: its assets are the 153 free monthly JKP factors of Jensen, Kelly and Pedersen, each described each month by its past 1-, 3- and 12-month returns, its 12-month volatility and its 36-month average return. Every model is re-estimated monthly on the previous 60 months and judged on the next month; Part D needs PyTorch, and the other parts run on a laptop in minutes.

\begin{itemize}
  \item \textbf{Part A, simple baselines.} The SDF from the Fama--French five factors plus momentum, and the shrinkage SDF of Equation~\eqref{eq:22-kns} on all 153 factors.
  \item \textbf{Part B, no attention.} Each factor's weight depends only on its own characteristics, linearly and through a small neural network.
  \item \textbf{Part C, linear attention.} The linear transformer of Equation~\eqref{eq:22-linear}, estimated as a closed-form ridge regression.
  \item \textbf{Part D, a small transformer.} One to four blocks, trained to maximize the SDF's Sharpe ratio and averaged over ten random starts.
  \item \textbf{Part E, scoring.} Out-of-sample Sharpe ratios, alphas against every other model, and pricing errors, for the full period and after 2002 (Lab Figures~\ref{ch:22}.1 to~\ref{ch:22}.3).
\end{itemize}

% [Draft note: the figures and numbers are generated by running the lab on data at the time of publication. The book prints the results; the text above does not anticipate them.]

\label{sec:22-stress}
\textbf{Robustness checks.} As a placebo, the lab shuffles characteristics across factors each month; if attention still adds alpha, it is fitting noise. Part E is repeated with factors built from large stocks only and with an assumed cost for each change in factor weights. Finally, following Nagel's (2025) test in Chapter~\ref{ch:20}, the lab compares the transformer with a simple rule that overweights factors with high recent returns and low volatility. If the rule explains the transformer's returns, attention has rediscovered a simple strategy. The lab ends by asking you to reach your own verdict on the scale of Section~\ref{sec:01-method} and compare it with Section~\ref{sec:22-verdict}.

\section{Where the evidence stands}
\label{sec:22-verdict}

The published studies point to an early picture: \textbf{one careful study found that a transformer prices U.S. stocks better than factor models and earlier machine-learning models out of sample. Its value after costs and among large stocks has not been tested, and no one has yet reproduced the result independently.} The table summarizes the evidence as of November 2026.

{\small
\begin{longtable}{@{}>{\raggedright\arraybackslash}p{0.19\linewidth}>{\raggedright\arraybackslash}p{0.13\linewidth}>{\raggedright\arraybackslash}p{\dimexpr 0.45\linewidth-6\tabcolsep\relax}>{\raggedright\arraybackslash}p{0.23\linewidth}@{}}
\toprule
Situation & Holds? & Why & Evidence \\
\midrule
\endfirsthead
\multicolumn{4}{@{}l}{\emph{(continued)}} \\
\toprule
Situation & Holds? & Why & Evidence \\
\midrule
\endhead
\midrule
\multicolumn{4}{r@{}}{\emph{(continued on the next page)}} \\
\endfoot
\bottomrule
\endlastfoot
Against factor models & \textbf{Yes} & Far higher Sharpe ratio and smaller pricing errors than the Fama--French six and $q$-factors. & Kelly, Kuznetsov, Malamud and Xu (2025) \\*
\addlinespace[3pt]
Against networks without attention & \textbf{Yes} & Pricing errors about 30\% smaller, with an alpha of about 1.9\% a year ($t = 8.4$). & Kelly, Kuznetsov, Malamud and Xu (2025) \\
\addlinespace[3pt]
Deep SDFs in general & \textbf{Yes} & A network trained on the pricing condition beat its benchmarks, mostly through familiar characteristics. & Chen, Pelger and Zhu (2024) \\
\addlinespace[3pt]
After 2002 & \textbf{Partly} & Smaller but still present: a Sharpe ratio of 3.37 against 3.07. & Kelly, Kuznetsov, Malamud and Xu (2025) \\
\addlinespace[3pt]
Large stocks & \textbf{Weakly} & Best within each size group, but much weaker among mega and large stocks. & Kelly, Kuznetsov, Malamud and Xu (2025) \\
\addlinespace[3pt]
After trading costs & \textbf{Untested} & No costs were reported. & None yet \\*
\addlinespace[3pt]
Independent replication & \textbf{Untested} & A working paper by a single team. & None yet \\
\end{longtable}
\addtocounter{table}{-1}% uncaptioned longtable: keep table numbers unchanged
}

\textbf{What would change the picture.} The economic form would gain support if a cost-aware version kept its advantage among large stocks and an independent team reproduced it. It would lose support if a simple factor-timing rule explained the transformer's returns.

\label{sec:22-open}
\textbf{Open questions.} What does attention learn, given that observable links between firms explain only a small part of it? And if a Sharpe ratio above four reflects mispricing rather than risk, how long can it last once the model is widely used?

\section{What practitioners and investors can use}
\label{sec:22-investors}

Read Sharpe ratios above three as research results, not as returns: they are hypothetical, frictionless and long--short, include hard-to-trade stocks, and no fund has delivered them for long. The lasting lesson is about information. Stocks are not independent, and an investor can use that idea without a transformer, for example by watching a company's suppliers, customers and competitors. Treat ``AI-powered'' as a claim to test, and ask a manager for the Sharpe ratio after costs among stocks the fund can trade, and how the model compares with a simple factor-timing rule.

\begin{callout}
\textbf{What Failed: the quant quake of August 2007.} In the first two weeks of August 2007, many quantitative equity funds suffered large, sudden losses, followed by a partial rebound. They held long--short portfolios built from similar characteristics such as value, momentum and quality, and Khandani and Lo (2011) traced the episode to the rapid sale of one or more large portfolios. The models were not wrong about average returns but about risk: funds that thought they held different strategies held nearly the same portfolio. A model that learns from the same public characteristics as everyone else may end up in the same crowded trades.
\end{callout}

% Key terms of the chapter, in alphabetical order, for readers new to finance.
% A term defined in an earlier chapter points to the chapter, section and page.
\providecolor{FEMaroon}{RGB}{128,30,45}
\providecolor{FECream}{RGB}{253,248,238}
\providecolor{FEGold}{RGB}{196,150,60}
\makeatletter
\@ifundefined{keyterms}{%
  \newtcolorbox{keyterms}[1][Key terms in this chapter]{breakable, enhanced,
    colback=FECream, colframe=FEGold, boxrule=0.6pt, arc=2pt, leftrule=3pt,
    left=6pt, right=6pt, top=3pt, bottom=3pt, before skip=8pt, after skip=8pt,
    fonttitle=\bfseries, coltitle=white, colbacktitle=FEMaroon, title={#1}}%
}{}
\makeatother

\begin{keyterms}
\footnotesize
\begin{description}[leftmargin=0pt, itemsep=0pt, parsep=0pt, topsep=0pt, font=\normalfont\bfseries\color{FEMaroon}]
  \item[Adversarial training.] Adversarial training pits two networks against each other, one searching for the portfolios that the other prices worst.
  \item[Attention.] Attention is the mechanism in a transformer that lets each input, such as a word or a stock, look at the other inputs to decide what it means.
  \item[Hansen--Jagannathan distance.] The Hansen--Jagannathan distance summarizes how far a model's SDF is from correctly pricing a set of test assets.
  \item[Linear transformer.] A linear transformer is a simplified transformer whose attention drops the exponential and the normalization.
  \item[Principal components.] See Chapter~\ref{ch:20}, Section~\ref{sec:20-models}, page~\pageref{sec:20-models}.
  \item[Query, key and value.] In attention, each input is turned into a query, a key and a value, and each input's query is compared with the others' keys to decide how much of their values to use.
  \item[Stochastic discount factor (SDF).] The stochastic discount factor is a single random variable that prices every asset, valuing payoffs more highly in bad times.
  \item[Transformer.] See Chapter~\ref{ch:21}, Section~\ref{sec:21-fm}, page~\pageref{kt:transformer}.
\end{description}
\end{keyterms}
